\section{Polynomial models with nonuniform contraction}
\label{sec:dynamics}

We retain the volume clock $s=\rho(o,t)$, its logarithmic-time speed
$q(s)=ds/d\tau=tR(o,t)$, and the charts of
Proposition~\ref{prop:shrinking-charts}.
Polynomial normalization and inverse iteration are classical tools
for attracting basins, notably in Rosay--Rudin's work
\cite[Appendix]{RosayRudin1988}; Chau--Tam brought this strategy into
K\"ahler uniformization \cite{ChauTamComplex2006,ChauTamStein,ChauTamSimply}.
The estimates proved here address the nonuniform contraction permitted
by Proposition~\ref{prop:ricci-windows}.  For sequences of holomorphic
automorphisms of $\C^n$, $n\ge2$, fixing zero and satisfying
$A|z|\le|f_j(z)|\le B|z|$ on one fixed ball about zero, with
$0<A<B<1$ independent of $j$, Bera--Verma prove that the attracting
basin is biholomorphic to $\C^n$ \cite[Theorem~1.1]{BeraVerma2024}.
The chart transitions here require the variable-rate estimates below.
In particular, we use only
\begin{equation}
 q\text{ nondecreasing},\qquad q\ge c>0,\qquad
 \liminf_{s\to\infty}\frac{q(s)}s\le2,\qquad
 \lambda_{\min}(t\Ric)(o)\ge b_0q(a_0s),
 \label{eq:dyn:rate-input}
\end{equation}
where $0<a_0<1$ and $b_0>0$ are fixed.

We continue to use $J_d$ and the fixed coefficient norms introduced in
Section~\ref{sec:charts}.
All maps and jets in this Section fix zero.  A coefficient bound
$\exp(o(s_k))$ means an upper bound $\exp(\eta_k s_k)$ with
$\eta_k\to0$; its rate of convergence may depend on parameters already
fixed.  No assertion of uniformity as those parameters vary is intended.
For polynomial automorphisms put
$G_{k,l}=G_{l-1}\circ\cdots\circ G_k$ and $G_{k,k}=\Id$.

\begin{proposition}[Polynomial models]\label{prop:polynomial-models}
Fix an integer $d\ge2$ and a macro width $0<h<1$.  After starting at a
sufficiently late macro endpoint, there is a grid $s_k\to\infty$ with
macro endpoints $S_j=e^{hj}$ and the following properties.  Each macro
interval is either one bad step or is divided into regular steps of
length comparable to $\sqrt{S_j}$.  Thus, eventually,
\begin{equation}
 c_h\sqrt{s_k}\le\Delta s_k:=s_{k+1}-s_k
 \le(e^h-1)s_k.
 \label{eq:dyn:grid-size}
\end{equation}
For $F_k=F_{s_k,s_{k+1}}$ there are polynomial local coordinate changes
$H_k$ and polynomial automorphisms $G_k$ of $\C^n$ satisfying
\begin{equation}
 J_d(H_{k+1}\circ F_k)=J_d(G_k\circ H_k).
 \label{eq:dyn:jet-conjugacy}
\end{equation}
The degrees of $H_k,G_k,G_k^{-1}$ are bounded independently of $k$.
The coefficients of $H_k$ and its inverse jet are $\exp(o(s_k))$;
$H_k'(0)$ and its inverse have polynomial bounds in $s_k$.
For a constant $C_d$ depending only on the fixed jet order and dimension,
\begin{align}
 \norm{G_k}_{\mathrm{coeff}}+\norm{G_k^{-1}}_{\mathrm{coeff}}
 &\le \exp\bigl(C_d\Delta s_k+o(s_k)\bigr),
 \label{eq:dyn:coefficient-bound}\\
 \norm{G_k'(0)^{-1}}
 &\le s_k^{C}\exp(\Delta s_k/2).
 \label{eq:dyn:center-inverse}
\end{align}
In particular, the coefficient $C_d$ of $\Delta s_k$ is independent of
$h$.  Finally, for some finite $D$,
\begin{equation}
 \deg G_{0,k}\le D s_k
 \quad\text{on an unbounded subsequence of macro endpoints.}
 \label{eq:dyn:forward-degree}
\end{equation}
\end{proposition}

First we normalize the Jacobian jets coherently along the grid.
On regular steps we then construct polynomial models whose nonlinear
terms and intervening linear changes preserve bounds for the degree of
each component.  On bad steps, realization of the full normalized jet
costs a fixed degree factor.  The last part of the proof charges these
factors to growth of $q$ and obtains \eqref{eq:dyn:forward-degree}.
The distinction between the constant $C_d$ in
\eqref{eq:dyn:coefficient-bound} and constants hidden in $o(s_k)$ will
be used in Section~\ref{sec:uniformization}.

\subsection{Coherent volume jets}

On the uniform transition ball, the germs $F_k$ are nonsingular and have
uniformly bounded derivatives of each fixed order on a smaller ball.
Their center derivatives are contractions, and chart normalization gives
\begin{equation}
 \abs{\det F_k'(0)}=e^{-\Delta s_k/2},\qquad
 \norm{F_k'(0)^{-1}}\le e^{\Delta s_k/2}.
 \label{eq:dyn:transition-linear}
\end{equation}
The second inequality follows because every singular value is at most
one and their product is $e^{-\Delta s_k/2}$.
Choose the holomorphic logarithm vanishing at zero of
\[
 P_k(z)=\log\frac{\det F_k'(z)}{\det F_k'(0)}.
\]
The upper bound on its real part is $C+\Delta s_k/2$.  The coefficient
estimate for a holomorphic function with bounded real part, applied on
concentric balls and with $P_k(0)=0$, therefore gives
$\norm{J_{d-1}P_k}\le C_d(1+\Delta s_k)$.

The volume correction uses only the spacing bounds
\eqref{eq:dyn:grid-size}, so it can be constructed before the macro
intervals are classified.  On any such grid the jet series
\begin{equation}
 \ell_k=\sum_{l\ge k}J_{d-1}
       \bigl(P_l\circ F_{s_k,s_l}\bigr)
 \label{eq:dyn:volume-series}
\end{equation}
converges.  Indeed Proposition~\ref{prop:shrinking-charts} bounds each
coefficient of the summand by
\[
 C_d(1+\Delta s_l)
       \min\{1,\exp(C_d s_k-\sigma s_l)\}.
\]
The absence of a constant term in $P_l$ is essential here.  For
$s_l\le C'_d s_k$ the step lengths telescope and the number of steps is
polynomial in $s_k$.  Beyond that range the displayed exponential is
summable, since the grid eventually has spacing at least one and
$\Delta s_l\le(e^h-1)s_l$.  Thus $\ell_k$ has polynomial coefficient
bounds in $s_k$.  The composition identity for transitions yields
\begin{equation}
 \ell_k=J_{d-1}\bigl(P_k+\ell_{k+1}\circ F_k\bigr).
 \label{eq:dyn:volume-cocycle}
\end{equation}

Let $E_k=J_{d-1}\exp(\ell_k)$ and define a polynomial change
\[
 A_k(z)=\left(z_1,\ldots,z_{n-1},
       \int_0^{z_n}E_k(z_1,\ldots,z_{n-1},\zeta)\,d\zeta\right).
\]
It is tangent to the identity, has degree at most $d$, and satisfies
$J_{d-1}\log\det A_k'=\ell_k$.  Its coefficients and inverse jet have
polynomial bounds in $s_k$.  In the new coordinates the normalized
logarithmic Jacobian jet of $A_{k+1}F_kA_k^{-1}$ is
\[
 J_{d-1}\bigl(\ell_{k+1}\circ F_k+P_k-\ell_k\bigr)\circ A_k^{-1}=0.
\]
This records the sign of the volume correction explicitly.  We shall
call a jet with Jacobian determinant equal, through degree $d-1$, to
its center determinant a \emph{constant-volume jet}.

\begin{lemma}[Realization of constant-volume jets]\label{lem:volume-jets}
For fixed $n,d$, every invertible constant-volume $d$-jet fixing zero
is the $d$-jet of a polynomial automorphism of $\C^n$.  Its degree and
inverse degree are bounded by a constant $D_d$.  Its coefficients and
inverse coefficients are bounded polynomially in the coefficients of
the given jet and its inverse linear part.  These bounds do not depend
on the grid or on $h$.
\end{lemma}

\begin{proof}
Normalize the derivative to the identity.  If a volume-preserving
polynomial automorphism agrees with the required jet through degree
$m-1$, the first discrepancy is a homogeneous vector polynomial $X_m$
of degree $m$ with $\operatorname{div}X_m=0$: this is the degree $m-1$
part of the Jacobian determinant condition.

The space of such fields is spanned by
\begin{equation}
 X(z)=l(z)^m v,\qquad l\in(\C^n)^*,\quad v\in\C^n,\quad l(v)=0.
 \label{eq:dyn:shear-fields}
\end{equation}
This is the divergence-free shear lemma of Anders\'en
\cite{AndersenVolume1990}; see also
\cite[proof of Theorem~3.1]{ForstnericLarusson}.
We include the finite-dimensional proof to track the coefficient bounds.
Pair a linear functional on
homogeneous vector polynomials with $l^m v$.  The result has the form
$A(l)(v)$, where $A(l)$ is a covector-valued homogeneous polynomial of
degree $m$.  If the functional annihilates all fields in
\eqref{eq:dyn:shear-fields}, then $A(l)$ is proportional to $l$.
The polynomial identities $l_iA_j(l)=l_jA_i(l)$ show that
$A(l)=B(l)l$ for a homogeneous polynomial $B$ of degree $m-1$.
Such functionals are exactly the image of the transpose of divergence,
since $\operatorname{div}(l^m v)=m l(v)l^{m-1}$.  Taking annihilators
proves the spanning assertion.  In dimension one the divergence kernel
in these degrees is zero, so this argument also covers $n=1$.

Select a finite spanning list of the fields
\eqref{eq:dyn:shear-fields}, once for each $2\le m\le d$.
The time-$c$ map of such a field is the polynomial shear
$z\mapsto z+c l(z)^m v$, whose inverse has the minus sign; both have
Jacobian determinant one.  A fixed linear decomposition of $X_m$
therefore gives a finite composition of shears correcting degree $m$
without changing lower degrees.  Induction through $m=d$, followed by
the original linear map, proves realization.  There are a fixed number
of shears, and each operation on the truncated jets is polynomial in
the preceding coefficients and the inverse linear part.  This proves
all the coefficient and degree bounds, including those for the inverse.
Restoring the original linear map gives the realizing automorphism
the determinant of that linear map; the correcting shears have
determinant one.
\end{proof}

\subsection{Rate order without commuting matrices}

Consider one interval in logarithmic time, of length $T>0$.
Singular-value decompositions in its start and end unitary tangent
coordinates give derivative
\begin{equation}
 D=\diag(e^{-\mu_1T/2},\ldots,e^{-\mu_nT/2}),\qquad
 0\le\mu_1\le\cdots\le\mu_n,\qquad
 \sum_i\mu_i=\frac{\Delta s}{T}.
 \label{eq:dyn:diagonal-rates}
\end{equation}
The volume changes $A_k$ do not change these linear parts.

We need an order statement for the rate forms, and not merely their
ordered eigenvalues.  Normalize the cometric at the start to $I$ in a
fixed covector basis.  In the diagonalized basis its value at the end
is $\diag(e^{\mu_iT})$.  On the complex strip $0\le\Re z\le T$, the
holomorphic covector with components $a_i e^{-\mu_i z/2}$ has squared
norm $\abs a^2$ on both boundary lines.  Its squared norm is bounded
and subharmonic by Theorem~\ref{thm:joint-positivity}, applied with
$t=t_{\mathrm{start}}e^{\Re z}$; equivalently the complex time variable
there is shifted by $z/2$.  The bounded subharmonic maximum principle
on the strip gives the same upper bound inside.  Consequently the
cometric is bounded above by the flat exponential interpolation of its
endpoint values.  Differentiating this matrix inequality at the two
ends gives
\begin{equation}
 (t\Ric)_{\mathrm{start}}\le\diag(\mu_i),\qquad
 \diag(\mu_i)\le(t\Ric)_{\mathrm{end}},
 \label{eq:dyn:endpoint-rate-order}
\end{equation}
where each matrix is expressed in the corresponding endpoint unitary
covector coordinates.  The inequalities have opposite endpoint signs
because the matrix difference vanishes at both ends and is
nonnegative inside.

At a shared endpoint of two steps, let $U_k$ send old tangent
components to new tangent components.  It is unitary; a row covector
$v$ in the new coordinates has old row $vU_k$.  Combining the two
inequalities in \eqref{eq:dyn:endpoint-rate-order} gives
\begin{equation}
 \sum_l\mu_l^{\mathrm{old}}\abs{(vU_k)_l}^2
 \le\sum_i\mu_i^{\mathrm{new}}\abs{v_i}^2.
 \label{eq:dyn:loewner-jump}
\end{equation}
This proof allows the instantaneous Ricci matrices to fail to commute.
In particular, the sorted rates are nondecreasing from step to step.

\subsection{The macro grid and propagated weights}

The positive weights below will bound ordinary polynomial degrees
component by component.  For a polynomial map
$P=(P_1,\ldots,P_n):\C^n\to\C^n$ with $\deg P_l\le w_l$, one has
\[
 \deg(P_1^{\beta_1}\cdots P_n^{\beta_n})
 \le\sum_l\beta_lw_l.
\]
Thus a nonlinear map preserves these degree bounds if every monomial
$z^\beta$ in component $i$ satisfies
$\sum_l\beta_lw_l\le w_i$.  A linear change between two weight systems
preserves the bounds provided every nonzero entry sends an old
coordinate to a new coordinate of at least its weight.  We shall
enforce both rules on regular steps.  Bad steps may multiply degrees
by a fixed factor; their defining rate growth will pay for that factor.

Put $Q_j=q(S_j)$ and choose an integer $N\ge1$ with
$e^{-hN}\le a_0$.  A macro index $j$ is \emph{flagged} when
\begin{equation}
 \log Q_{j+1}-\log Q_{j-N}>M,
 \label{eq:dyn:flag}
\end{equation}
where the fixed threshold $M$ will be chosen sufficiently large below.
Declare bad the union of the expanded flag intervals
$[j-N,j+1+N)$ in the integer macro index line.  Each bad macro is a
single step.  Divide every remaining macro interval into equal steps
whose lengths are comparable to $\sqrt{S_j}$, for example by taking
the number of pieces to be the nearest larger integer to
$(S_{j+1}-S_j)/\sqrt{S_j}$.  Once $j$ is large this gives
\eqref{eq:dyn:grid-size}.

If macro $j$ is regular, it is not flagged.  On each of its steps,
\eqref{eq:dyn:rate-input} and \eqref{eq:dyn:endpoint-rate-order} imply
\begin{equation}
 b_0Q_{j-N}\le\mu_i\le Q_{j+1},\qquad
 bQ_j\le\mu_i\le b^{-1}Q_j
 \label{eq:dyn:regular-rates}
\end{equation}
for a fixed $b>0$ depending on $M,b_0$.  These constants may depend on
$h$ after all parameter choices have been made.

For a regular macro set $\theta_j=Q_j/j^2$ and choose bounded increasing
multipliers $\lambda_j$ on the whole macro sequence by
\[
 \lambda_{j+1}/\lambda_j=1+K_0/j^2.
\]
Their product is bounded, since $\sum j^{-2}<\infty$; an overall
constant factor is free.  Choose $0<\epsilon_0<1/(10n)$.
For each regular macro, form a finite directed graph whose vertices
are the pairs $(k,i)$ consisting of a step in that macro and one of
its coordinates.  Give this vertex raw weight
$\lambda_j(\mu_i-\theta_j)$, where $\mu_i$ is its rate in step $k$.
Draw an edge to a vertex in the same or the next step precisely when
the destination rate is at least the source rate minus
$\epsilon_0\theta_j$.  Define $w_{k,i}$ as the maximum raw weight
among vertices from which $(k,i)$ is reachable, allowing paths of
length zero.  This maximum exists because the graph is finite, even
when same-step edges form cycles.  Assign the weight to both ends of
its diagonal step.

This propagation loses at most $n\epsilon_0\theta_j$ in rates,
independently of the number of steps.  To verify the assertion, follow
any permitted chain and keep its successive lower records.  Once a
sorted index has appeared, its rate at a later occurrence is at least
its earlier rate, by \eqref{eq:dyn:loewner-jump}.  A new lower record
can therefore occur only on the first visit to an index; each such
record lowers the previous record by at most
$\epsilon_0\theta_j$.  There are at most $n$ indices.  It follows that
\begin{equation}
 \lambda_j(\mu_i-\theta_j)\le w_{k,i}
 \le\lambda_j(\mu_i-0.9\theta_j).
 \label{eq:dyn:weights}
\end{equation}
In particular, whenever a permitted passage is possible the weights
do not decrease.

The same nondecrease can be required across adjacent regular macros.
Indeed, write $\mu'$ for a new rate and $\mu$ for an old one, with
$\mu'\ge\mu-\epsilon_0\theta_j$.  A lower bound on the new raw weight
minus the upper bound on the old propagated weight is
\[
 (\lambda_{j+1}-\lambda_j)\mu'
 -\lambda_{j+1}\theta_{j+1}
 +(0.9-\epsilon_0)\lambda_j\theta_j.
\]
Since $\mu'\ge bQ_{j+1}$, choosing $K_0$ large in terms of $b$ makes
this expression nonnegative at every sufficiently large index.
Thus the passage rule is valid also at regular macro junctions.
Increase the overall scale of $\lambda$ so that all weights are at
least one after the eventual starting index.

Call a nonlinear monomial $z^\beta$ in component $i$ \emph{allowed}
when $\sum_l\beta_lw_{k,l}\le w_{k,i}$, and \emph{forbidden} otherwise.
For a forbidden monomial of degree $m\ge2$,
\eqref{eq:dyn:weights} gives
\begin{equation}
 \sum_l\beta_l\mu_l-\mu_i
 >(0.9m-1)\theta_j\ge0.8\theta_j.
 \label{eq:dyn:forbidden-gap}
\end{equation}
For an allowed monomial, every variable occurring in it has strictly
smaller weight than its output: the weights are positive and $m\ge2$.
Thus allowed nonlinear maps are triangular after ordering coordinates
by weight.

\subsection{Aligning the linear jumps}

\begin{lemma}[Linear alignment]\label{lem:linear-alignment}
On every run of regular steps there are unipotent matrices $L_k^+$,
with off-diagonal entries only in rows of smaller weight than their
columns, such that, on putting $L_k^-=D_k^{-1}L_k^+D_k$, the diagonal
step derivative remains $D_k$ and every regular jump becomes
\begin{equation}
 J_k=L_{k+1}^-U_k(L_k^+)^{-1},\qquad
 (J_k)_{il}=0\quad\hbox{if }w_{k+1,i}<w_{k,l}.
 \label{eq:dyn:aligned-jump}
\end{equation}
The matrices $L_k^+$ and their inverses have polynomial bounds in the
macro index $j$, while
\begin{equation}
 \norm{L_k^- -I}
 \le j^C\exp\left(-c\frac{\sqrt{S_j}}{j^2}\right).
 \label{eq:dyn:small-start-change}
\end{equation}
The estimates are uniform in finite terminal horizons, and hold on an
infinite regular run by a coefficientwise limit.
\end{lemma}

\begin{proof}
Work backward, taking $L_k^+=I$ at the last step of a finite run or
horizon.  Suppose the next small start adjustment is known and set
$B=L_{k+1}^-U_k$.  For a real threshold $x$, let $A_x$ be the span of
rows of $B$ whose new weights are at most $x$, and let $E_x$ be the old
coordinate row space with weights at most $x$.

Projection $P_{E_x}:A_x\to E_x$ is injective, and its inverse on its
image has norm $O(j)$.  For the proof assume $A_x\ne0$ and let $r_x$
be the largest new rate among its defining rows.  Every old rate
outside $E_x$ exceeds $r_x+\epsilon_0\theta_j$: otherwise the passage
rule would make that old weight at most the weight of a selected new
row.  For any unit $a\in A_x$, \eqref{eq:dyn:loewner-jump} and the
small next adjustment imply that its old rate Rayleigh quotient is at
most $r_x+o(\theta_j)$.  To see the size of this error, express $a$ as
a combination of the selected rows of $L_{k+1}^-U_k$; its coefficient
norm is $1+o(1)$, and the error is bounded by the operator norm of
$L_{k+1}^- -I$ times the largest old or new rate.  The regular rate
bounds make that error $o(\theta_j)$, also at a macro junction.
Nonnegativity of all old rates now gives
\[
 (r_x+\epsilon_0\theta_j)
       (1-\norm{P_{E_x}a}^2)
 \le r_x+o(\theta_j).
\]
As $r_x\le C Q_j$, this implies
$\norm{P_{E_x}a}^2\ge c/j^2$.  Empty or full coordinate spaces give
the same assertion in its evident form.

We give the graph extension explicitly.  List the distinct spaces
$E_x$ increasingly.  For each such space retain the largest required
$A_x$ with that $E_x$, so that the required $A_x$ remain nested.
Starting with the full row space, descend through the list.  Suppose
the next larger space $V_+$ already projects isomorphically to $E_+$
and contains the current $A=A_x$.  Write $T_+:E_+\to V_+$ for its
inverse projection, $E=E_x$, $Y=P_EA$, and $R_A:Y\to A$ for the
inverse of $P_E|_A$.  Define
\begin{equation}
 T_E(e)=R_A(P_Ye)+T_+((I-P_Y)e),\qquad e\in E.
 \label{eq:dyn:graph-extension}
\end{equation}
Here orthogonal projection $P_Y$ is taken in $E$.  Both terms lie in
$V_+$, projection to $E$ is $e$, and the range contains $A$.
Moreover $\norm{T_E}\le\norm{R_A}+\norm{T_+}$.
There are at most $n$ nontrivial stages, so the resulting graph lifts
have polynomial bounds in $j$.

For each coordinate row, take its lift at its own weight threshold.
These rows form $L_k^+$.  The rows up to every threshold span the
corresponding nested graph: they lie there and their projections form
a triangular basis of its coordinate space.  Each lifted row equals
its coordinate row plus terms of strictly higher weight.  Hence
$L_k^+=I+N_k$ with $N_k$ nilpotent, and
$(L_k^+)^{-1}=\sum_{r=0}^{n-1}(-N_k)^r$ has polynomial bounds.
The inclusion $A_x(L_k^+)^{-1}\subset E_x$ is exactly the zero pattern
in \eqref{eq:dyn:aligned-jump}.

For every nonzero off-diagonal entry in row $i$, column $l$, the
within-step passage rule gives $\mu_l-\mu_i>\epsilon_0\theta_j$.
Conjugation therefore gives
\[
 (D_k^{-1}L_k^+D_k)_{il}
 =(L_k^+)_{il}\exp\bigl(- (\mu_l-\mu_i)\Delta\tau_k/2\bigr).
\]
By \eqref{eq:dyn:regular-rates},
$\theta_j\Delta\tau_k\ge c\sqrt{S_j}/j^2$.
This proves \eqref{eq:dyn:small-start-change}; its exponential absorbs
every polynomial bound needed in the preceding backward step.
After a sufficiently late start the induction thus closes uniformly
in the horizon.  A diagonal subsequence of finite-horizon matrices
proves the assertion on infinite runs; the zero conditions and all
bounds are closed under that limit.
\end{proof}

\subsection{Homological equations and polynomial automorphisms}

Choose a single primary coordinate system at each grid endpoint.
Use the adjusted start coordinates if a regular step starts there;
otherwise use the adjusted end coordinates of a regular step just
ended, if there is one; use the volume coordinates in all other cases.
These choices are a linear change, of polynomial size together with
its inverse, following $A_k$.  For a regular step, before the possible
jump to the next regular start, write the resulting transition as
$\widehat F_k$, with derivative $D_k$.  Its actual transition to the
next primary coordinates is $J_k\widehat F_k$, with $J_k=I$ when the
next step is bad.  It has a constant-volume $d$-jet and polynomially
bounded coefficients.  This follows by formal composition and
inversion of the primary coordinate jets; only their local germs are
used.

We construct tangent-to-identity polynomial jets $h_k$ at regular
starts, taking $h_k=I$ elsewhere.  For a regular step seek
$G_k=J_kg_k$, where $g_k$ has linear part $D_k$ and only allowed
nonlinear terms, satisfying
\begin{equation}
 g_k\circ h_k
 =J_k^{-1}\circ h_{k+1}\circ J_k\circ\widehat F_k
 \quad\text{through degree }d.
 \label{eq:dyn:homological-identity}
\end{equation}
In homogeneous degree $m$, once lower degrees have been fixed, this
is
\begin{equation}
 g_k^{(m)}+D_kh_k^{(m)}
 =J_k^{-1}h_{k+1}^{(m)}(J_kD_kz)+B_k^{(m)}(z),
 \label{eq:dyn:homological-layer}
\end{equation}
where $B_k^{(m)}$ contains only already known data.
Put only forbidden coefficients into $h_k^{(m)}$ and only allowed
coefficients into $g_k^{(m)}$.  If $\Pi_k$ projects to the forbidden
coefficients, the equation for $h_k^{(m)}$ is a backward affine
recursion with linear operator
\begin{equation}
 T_kX=\Pi_k\left[D_k^{-1}J_k^{-1}X(J_kD_kz)\right].
 \label{eq:dyn:backward-operator}
\end{equation}
Conjugation by $J_k$ has polynomial norm in $j$ at each fixed degree.
After it has been expanded, a coefficient in output $i$ and input
$\beta$ receives the diagonal factor
$\exp[-(\sum_l\beta_l\mu_l-\mu_i)\Delta\tau_k/2]$.
Projection and \eqref{eq:dyn:forbidden-gap} consequently give
\begin{equation}
 \norm{T_k}\le j^{C_m}
       \exp\left(-c\frac{\sqrt{S_j}}{j^2}\right).
 \label{eq:dyn:backward-contraction}
\end{equation}
This estimate holds for every incoming homogeneous polynomial $X$;
no support restriction on $X$ has been imposed before conjugation.

Solve first on a finite regular horizon with terminal correction the
identity, and proceed by degree.  Inductively the lower-layer forcing,
including division by $D_k$, has coefficients bounded by
$\exp(C'_m\sqrt{s_k})$: regular steps have
$\Delta s_k\asymp\sqrt{s_k}$, and
$\norm{D_k^{-1}}\le e^{\Delta s_k/2}$.
The envelopes $\exp(C_m\sqrt{s_k})$ have bounded consecutive ratios,
because $\sqrt{s_{k+1}}-\sqrt{s_k}$ is bounded on regular steps.
Thus \eqref{eq:dyn:backward-contraction} is eventually less than
one half after weighting by these envelopes.  The backward geometric
series solves \eqref{eq:dyn:homological-layer} with
\begin{equation}
 \norm{h_k}+\norm{J_dh_k^{-1}}+\norm{g_k}
       +\norm{g_k^{-1}}_{\mathrm{coeff}}
 \le\exp(C_d'\sqrt{s_k}),
 \label{eq:dyn:regular-coefficients}
\end{equation}
uniformly in the terminal horizon.  Here $C_d'$ may depend on all
fixed grid parameters.  Taking coefficientwise diagonal limits gives
the same construction on an infinite regular run.

Each $g_k$ is a triangular polynomial automorphism: all nonlinear
inputs have smaller weights than their output.  Its inverse is found
by successive substitution in increasing weight order.  Its degree
and inverse degree are bounded in terms of $n,d$, and its determinant
is the constant $\det D_k$.

The coordinate jets $h_k$ also preserve volume through degree $d-1$.
This is not a property assumed of the forbidden projection.  In a
finite-horizon problem, differentiate the complete jet identity
\eqref{eq:dyn:homological-identity}.  The determinant of
$\widehat F_k'$ has the same constant jet as $\det D_k$, and $g_k$ has
that determinant exactly.  If $\det h_{k+1}'=1$ through degree $d-1$,
then the determinant identity forces the same statement for $h_k$.
Backward induction from the identity terminal jet proves it, and it
passes to the coefficientwise limits.

At a bad macro step realize, by Lemma~\ref{lem:volume-jets}, the jet of
the transition between the two primary coordinate systems after their
endpoint corrections $h_k,h_{k+1}$.  It is a constant-volume jet by
what was just proved.  Enlarge $D_d$ to bound the degrees and inverse
degrees of these models and of the regular models.
Let $H_k$ be the degree-$d$ truncation of the total primary change
followed by $h_k$.  It has the required center derivative and inverse
jet, and \eqref{eq:dyn:jet-conjugacy} follows.

For precision, on a bad step the realized input jet has coefficient
bound $\exp(o(s_k))$, whereas its inverse linear part is bounded by
$\exp(\Delta s_k/2+o(s_k))$, by
\eqref{eq:dyn:transition-linear}.  Lemma~\ref{lem:volume-jets} costs a
fixed polynomial in these two quantities.  Its polynomial exponent
depends only on $n,d$.  Therefore it gives
\eqref{eq:dyn:coefficient-bound} with a coefficient $C_d$ independent
of $h,N,M$.  All endpoint and regular costs
$\exp(C(d,h,N,M)\sqrt{s_k})$ remain in $\exp(o(s_k))$.
Finally the linear part of the jet conjugacy is
$G_k'(0)=H_{k+1}'(0)F_k'(0)H_k'(0)^{-1}$; the polynomial bounds on the
endpoint linear changes give \eqref{eq:dyn:center-inverse}.

\subsection{Charging bad components to rate growth}

\begin{lemma}[Forward degree control]\label{lem:degree-control}
The flag threshold $M$ can be chosen so that regular macro starts
occur arbitrarily late, and the polynomial models constructed above,
initialized at any sufficiently late such start, satisfy
\eqref{eq:dyn:forward-degree}.
\end{lemma}

\begin{proof}
Set
\[
 W=2N+1,\qquad D_*=D_d,\qquad
 B_* =\max\{0,\log(2/b_0)\},\qquad \alpha=\frac{M}{3W}.
\]
Choose $M$ so large that
\begin{equation}
 \alpha\ge\log D_*+B_*/W,\qquad \alpha>h.
 \label{eq:dyn:flag-threshold}
\end{equation}
A flag $f$ has expanded interval $I_f=[f-N,f+1+N)$, of length $W$,
and growth interval $[f-N,f+1]$, on which $\log Q$ increases by more
than $M$.

Consider a full finite bad component $[u,v)$.  Choose a maximal
disjoint collection of its expanded flag intervals.  Their concentric
triples cover the component: any unselected equal-length interval
meets a selected one and is contained in its triple.  Thus there are
at least $(v-u)/(3W)$ selected intervals.  Their growth intervals are
disjoint up to endpoints and lie in $[u,v-N]$.  Monotonicity gives
\begin{equation}
 \log\frac{Q_{v-N}}{Q_u}\ge\alpha(v-u).
 \label{eq:dyn:full-bad-growth}
\end{equation}
Every full component has length $L=v-u\ge W$, so
\eqref{eq:dyn:flag-threshold} implies
\begin{equation}
 D_*^L Q_u\le e^{-B_*L/W}Q_{v-N}
             \le(b_0/2)Q_{v-N}.
 \label{eq:dyn:bad-reset-cost}
\end{equation}

We now track the ordinary degree of each component of the accumulated
forward polynomial.  At a regular start bound these degrees by the
respective weights.  Allowed nonlinear terms preserve that bound
under composition, and \eqref{eq:dyn:aligned-jump} preserves it at a
regular jump.  If the next macro is bad, retain the last end weights;
at entry to $[u,v)$ their maximum is at most
$\lambda_{u-1}Q_u$.  Each bad map multiplies maximum degree by at most
$D_*$.  At the next regular macro $v$, its flag test fails, and hence
$Q_v\le e^M Q_{v-N}$.  For sufficiently large $v$ every new start
weight is at least
\[
 \lambda_v\bigl(b_0Q_{v-N}-Q_v/v^2\bigr)
 \ge\lambda_v(b_0/2)Q_{v-N}.
\]
As $\lambda_v\ge\lambda_{u-1}$,
\eqref{eq:dyn:bad-reset-cost} resets all degree bounds exactly to the
new weights.  There is no multiplicative loss for a completed bad
component.

Prefixes require a separate estimate.  For a bad prefix $[u,j)$,
the expanded flag intervals wholly contained in the prefix cover
$[u,j-W)$, when this interval is nonempty.  Indeed an expanded interval
covering a point before $j-W$ ends before $j$, and none crosses the
left endpoint $u$ of the component.  The same disjoint selection gives
\begin{equation}
 \log(Q_j/Q_u)\ge\alpha(j-u-W)_+,
 \qquad D_*^{j-u}Q_u\le D_*^W Q_j.
 \label{eq:dyn:bad-prefix-growth}
\end{equation}
This proof also applies to a prefix of an infinite bad component.
Such a component would give
$\log(Q_j/S_j)\ge(\alpha-h)j-O(1)\to\infty$.
It is impossible: choose continuous-clock points $x_r\to\infty$ with
$q(x_r)\le3x_r$, and let $S_{j_r}\le x_r<S_{j_r+1}$.
Then $Q_{j_r}\le3e^h S_{j_r}$, contradicting that growth.
Thus there are regular macro starts arbitrarily far out.

Start at one of them, after every fixed-parameter threshold above,
and label this grid point $k=0$, retaining the original large macro
indices $j$.  The accumulated map is initially the identity and its
component degrees are one, so the initial weight bound holds after
the overall scaling already allowed for $\lambda$.  Regular
propagation, the exact full-component reset, and the single prefix
loss in \eqref{eq:dyn:bad-prefix-growth} give, at all macro endpoints,
\[
 \deg G_{0,k(j)}\le
 C\max\{e^M,D_*^W\}Q_j,
\]
where $C$ can be taken to depend on the bounded supremum of the
multipliers.  On a regular macro one uses
$Q_{j+1}\le e^M Q_j$; inside a bad component one uses the prefix
estimate.  Applying the bound at the endpoints $j_r$ just selected
proves \eqref{eq:dyn:forward-degree}.  This uses only a subsequence
with $Q_j=O(S_j)$, never an all-time bound on $q(s)/s$.
\end{proof}

Lemmas~\ref{lem:volume-jets}, \ref{lem:linear-alignment}, and
\ref{lem:degree-control}, together with the homological construction,
complete the proof of Proposition~\ref{prop:polynomial-models}.
